%ValueSemantics.tex

\ksec{Feature Values Semantics (KerML 8.4.4.11)}\label{sec:value}

\texttt{FeatureValue} either binds a value to a feature permanently (=), or assigns an initial value (:=).  Type checking must ensure the type of the expression is the type of the feature.

\begin{restatable}{definition}{blbfv}~\blTitle{df-bl.bfv}~ Define Bound Feature Value
\begin{align}
\texttt{classifier C \{ feature f:F=e; \} }~\Rightarrow \notag \\
\vdash \forall \tau \in \texttt{life}(C)~\vert~  \textbf{I}\langle \texttt{f},\tau\rangle~=~e
\label{eq:df-bl.bfv}\tag{df-bl.bfv}
\end{align}
\end{restatable}

\pn \texttt{df-bl.bfv} is book-original, with no \texttt{SFS.mm} formalization of its own (confirmed
absent from \texttt{SFS.mm} by grep) -- \texttt{=} binding a feature's value permanently is KerML's
own concrete syntax, not a domain predicate ported from a real \texttt{.kerml} file's \texttt{@Assert}.
It does have real Lean backing below.

\begin{restatable}{leandefinition}{blbfvLean}~\leanTitle{BoundFeatureValue}~ Define bound feature
value in Lean: \texttt{f}'s value is the fixed \texttt{e} throughout \texttt{C}'s whole
\texttt{life}, via \texttt{Get}.

\texttt{\textcolor{leancolor}{def} \textcolor{leanyellow}{BoundFeatureValue} (C : Occurrence) (f : Element) (e : Set Item) : Prop :=}

\texttt{~~$\forall$ tau : Time, tau.val $\in$ life C $\to$ Get C f tau = e}
\label{BoundFeatureValue}
\end{restatable}

\begin{restatable}{leantheorem}{dfBlBfvLean}~\leanTitle{df\_bl\_bfv}~ The book's own \ref{eq:df-bl.bfv}
biconditional, in Lean -- free directly from \texttt{BoundFeatureValue}'s own definition.

\texttt{\textcolor{leancolor}{theorem} \textcolor{leanyellow}{df\_bl\_bfv} (C : Occurrence) (f : Element) (e : Set Item) :}

\texttt{~~BoundFeatureValue C f e $\leftrightarrow$ ($\forall$ tau : Time, tau.val $\in$ life C $\to$ Get C f tau = e) := Iff.rfl}
\label{df_bl_bfv}
\end{restatable}

Setting an initial feature value only makes sense for variable features:  \texttt{in}, \texttt{out}, \texttt{inout}, and \texttt{var},

\begin{restatable}{definition}{blifv}~\blTitle{df-bl.ifv}~ Define Initial Feature Value
\begin{align}
\texttt{classifier C \{ var feature f:F := e; \} }~\Rightarrow \notag \\
\vdash (~\tau = \texttt{birth}(C)~\wedge~  \textbf{I}\langle \texttt{f},\tau\rangle~=~e ~)
\label{eq:df-bl.ifv}\tag{df-bl.ifv}
\end{align}
\end{restatable}

\pn \texttt{df-bl.ifv} is book-original, with no \texttt{SFS.mm} formalization of its own (confirmed
absent from \texttt{SFS.mm} by grep) -- \texttt{:=} assigning a feature's initial value is likewise
KerML's own concrete syntax. It does have real Lean backing below.

\begin{restatable}{leandefinition}{blifvLean}~\leanTitle{InitialFeatureValue}~ Define initial feature
value in Lean: \texttt{f}'s value is \texttt{e} specifically at \texttt{C}'s \texttt{birth}, via
\texttt{Get} (unlike \texttt{BoundFeatureValue}, only at one instant, not throughout \texttt{life}).

\texttt{\textcolor{leancolor}{def} \textcolor{leanyellow}{InitialFeatureValue} (C : Occurrence) (f : Element) (e : Set Item) : Prop :=}

\texttt{~~Get C f (birth C) = e}
\label{InitialFeatureValue}
\end{restatable}

\begin{restatable}{leantheorem}{dfBlIfvLean}~\leanTitle{df\_bl\_ifv}~ The book's own \ref{eq:df-bl.ifv}
biconditional, in Lean -- free directly from \texttt{InitialFeatureValue}'s own definition.

\texttt{\textcolor{leancolor}{theorem} \textcolor{leanyellow}{df\_bl\_ifv} (C : Occurrence) (f : Element) (e : Set Item) :}

\texttt{~~InitialFeatureValue C f e $\leftrightarrow$ Get C f (birth C) = e := Iff.rfl}
\label{df_bl_ifv}
\end{restatable}


